En este trabajo diseñamos la capa de datos de un copiloto RAG interno para una distribuidora mayorista. El problema que buscamos resolver es que la información relevante está repartida entre documentos, procedimientos, condiciones comerciales y datos operativos, por lo que encontrar una respuesta confiable puede ser lento y depender del conocimiento informal de las personas.

La solución propuesta usa PostgreSQL como motor principal, junto con las extensiones \texttt{pgvector} y \texttt{btree\_gist}. En la misma base se guardan los datos operativos, los metadatos de los documentos, los fragmentos, los embeddings, los permisos y los eventos de auditoría. Elegimos este enfoque porque permite aplicar la vigencia de los documentos y los permisos antes de hacer la búsqueda por similitud. Así evitamos recuperar una fuente vencida o que el usuario no tiene permitido ver.

La implementación incluye datos sintéticos reproducibles, nueve consultas representativas, restricciones de integridad, control de acceso mediante RLS y evidencias de las respuestas generadas. El trabajo no implementa una interfaz, una API ni un modelo de lenguaje real: el objetivo es demostrar que la capa de datos puede sostener ese tipo de aplicación de manera consistente, trazable y con posibilidades de crecimiento.
